\section{Development and Special Tasks (Week 2--5) - Fahmida}

\subsection{Outline of the design}

The protocol family is developed incrementally from a minimal delegation model toward a structured IDP-based authorization design.
The central objective is to allow user $A$ to delegate controlled read access at server $B$ to service $P$, without sharing long-term credentials.
Weeks 2--4 focus on formalizing the threat model, checking executability under the lazy intruder, and improving robustness through typed message formats and explicit key-distribution assumptions.

\subsection{Description of protocol development and OFMC findings}

\subsubsection{Week 2: Dolev--Yao and static analysis}

In Week 2 we used a simplified baseline protocol to perform the required static Dolev--Yao analysis under a passive intruder (observe-only, no message injection).
In this setting, the intruder learns public and visible fields (agent identities and request context).
The intruder does not derive private keys, and cannot decrypt ciphertext encrypted for an uncompromised honest recipient.
This baseline analysis gave us a clear starting point: some structures are too weak and must be refined in later weeks to avoid confidentiality/authentication issues in active settings.

\subsubsection{Week 3: Lazy intruder}

Week 3 introduces identity-provider-issued authorization and the special-task executability check.
We verified that role $P$ is executable when $P$ is instantiated by the intruder:
the intruder can construct each required outgoing message from previously available knowledge.
Operationally, the intruder can receive a token from $A$, forward it to $B$, and continue protocol interaction.
This is expected for a delegated service role.
The intended security condition is not ``intruder cannot send,'' but rather ``only correctly issued IDP-bound authorization should be accepted by $B$.'' 
OFMC confirms that this version is executable and also reports attacks on broader goals (\texttt{ATTACK\_FOUND}),
which is consistent with an intermediate development stage where functionality and threat-model realism are added before full hardening.

\subsubsection{Week 4: Type-flaw resistance and GetPublicKey}

In Week 4 we introduced explicit format constructors (e.g., request/token/data wrappers) to reduce message-type confusion.
The core typing claim is: for any two non-variable message patterns, if they unify, they belong to the same intended kind/type.
We also added a separate GetPublicKey protocol so that $A$, initially trusting only $pk(IDP)$, can obtain $pk(X)$ from the identity provider.
This removes the unrealistic assumption that $A$ initially knows every party public key.
In verification terms, this week improved model discipline but did not eliminate all attacks:
the key-lookup model remains executable and still produces a weak-auth finding in OFMC, which is documented as a design limitation rather than hidden.

\subsubsection{Week 5: Channels and pseudonymous TLS abstraction}

In Week 5 we replaced as much direct transport cryptography as possible with pseudonymous channels, following the lecture model for TLS-style communication.
The goal of this step was compositional clarity: keep the authorization logic explicit while abstracting channel guarantees.
Concretely, owner--IDP and service--server communications were modeled with channel notation instead of only low-level encryption syntax.

OFMC result for this variant is \texttt{ATTACK\_FOUND} on secrecy, with leaked term of the form \texttt{fData(Photos(...))}.
This outcome is technically important: channel abstraction can reduce protocol complexity, but it does not automatically fix authorization logic flaws.
Therefore, Week 5 is recorded as a valid and useful intermediate snapshot in the development chain.

\subsection{AnB formalization}

\subsubsection{Design decisions and development}

The early versions used stronger initial assumptions to get an executable baseline quickly.
During development, these assumptions were reduced:
\begin{itemize}
  \item direct trust assumptions were replaced by IDP-issued authorization,
  \item typed formats were added to improve message discipline,
  \item key lookup was split into a dedicated protocol before main delegation flow when needed.
\end{itemize}

This staged approach made attack traces and modeling limitations visible at each step and allowed targeted refinement.
In particular, Weeks 3--5 show a clear methodological progression: executable model $\rightarrow$ stronger structure/typing $\rightarrow$ channel abstraction, with each step validated by OFMC and discussed from a security-engineering perspective.

\subsubsection{Modeling in AnB}

\begin{lstlisting}
Protocol: OpenAuth_W4_KeyLookup

Types:
  Agent A,I,X;
  Function pk,fKeyReq,fKeyResp;

Knowledge:
  A: A,I,X,pk(A),pk(I),inv(pk(A)),fKeyReq,fKeyResp;
  I: A,I,X,pk(A),pk(I),pk(X),inv(pk(I)),fKeyReq,fKeyResp;

Actions:
  A -> I: fKeyReq(A,X)
  I -> A: {fKeyResp(X,pk(X))}inv(pk(I))

Goals:
  A authenticates I on fKeyResp(X,pk(X));
\end{lstlisting}

For Week 5, the corresponding channel-based model is:
\begin{lstlisting}
Protocol: OpenAuth_W5_PseudonymousChannels
...
Actions:
  [p] *->* A: p,B,F,T,R
  [A] *->* I: A,p,B,F,T,R,h(pw(A,I),R)
  I *->* [A]: fGrant(A,p,B,F,T,Grant)
  [A] *->* p: fGrant(A,p,B,F,T,Grant)
  [p] *->* B: fUse(fGrant(A,p,B,F,T,Grant),F)
  B *->* [p]: fData(Photos)
\end{lstlisting}

